\chapter{DataFrames}
\label{cap:dataframes}

\section{O que é um DataFrame}

Um DataFrame é uma tabela bidimensional com colunas nomeadas. Em \hayashi{}, cada
coluna é um vetor de ponto flutuante de 64 bits. Valores ausentes são representados
por \texttt{NaN} (\textit{Not a Number}).

O DataFrame é o tipo central da linguagem — praticamente toda análise econométrica
começa com um DataFrame carregado de alguma fonte.

\section{Semântica de cópia (COW)}

DataFrames usam \textit{copy-on-write} (COW): atribuir um DataFrame a outra variável
cria um alias que compartilha os mesmos dados internos. A cópia real ocorre apenas no
momento em que um dos dois é modificado.

\begin{lstlisting}[caption={COW: baseline preservado}]
load "painel.csv" as df

let baseline = df      // alias barato, sem cópia

let df2 = df |> generate(z = x^2)    // df2 tem 'z', df não
let df3 = df |> filter(ano > 2000)   // df3 filtrado, df intacto

// baseline == df original em qualquer ponto
\end{lstlisting}

Esta semântica permite criar variantes de análise sem custo de memória até que
uma modificação real seja necessária.

\section{Criando DataFrames}

\subsection*{Bloco \texttt{input}}

Para pequenos datasets inline — ideal para exemplos e testes:

\begin{lstlisting}[caption={input block}]
input x y z
  1   2  10
  3   4  20
  5   6  30
end
\end{lstlisting}

O bloco \lstinline{input} cria um DataFrame chamado \lstinline{df} por padrão.
O número de colunas é definido pelo cabeçalho; cada linha de dados deve ter o
mesmo número de valores.

\begin{nota}
  \lstinline{input} aceita apenas valores numéricos. Use \lstinline{.} para
  representar valores ausentes. Para dados com colunas de texto, use
  \lstinline{load}.
\end{nota}

\subsection*{\texttt{load} — arquivos externos}

\begin{lstlisting}[caption={Carregando arquivos}]
load "dados.csv" as vendas
load "painel.dta" as painel        // formato Stata
load "resultado.parquet" as df
load "planilha.xlsx" as df         // primeira aba
load "planilha.xlsx" sheet="Dados" as df
\end{lstlisting}

Formatos suportados: CSV, Parquet, Stata (\texttt{.dta}), Excel (\texttt{.xlsx}),
SQLite.

\section{Salvando DataFrames}

\begin{lstlisting}
save df "saida.csv"
save df "saida.parquet"
save df "saida.dta"
save df "saida.xlsx"
\end{lstlisting}

\section{Inspecionando DataFrames}

\subsection*{\texttt{display}}

Exibe as primeiras linhas do DataFrame:

\begin{lstlisting}
display df
display df, 20    // exibe até 20 linhas
\end{lstlisting}

\subsection*{\texttt{describe}}

Exibe a estrutura: número de observações, variáveis e tipos:

\begin{lstlisting}
describe df
\end{lstlisting}

Saída típica:

\begin{lstlisting}[style=inline]
DataFrame: 1250 obs, 8 vars
  ano     Float  (0 missing)
  uf      Float  (0 missing)
  renda   Float  (12 missing)
  gini    Float  (3 missing)
\end{lstlisting}

\subsection*{\texttt{count} e \texttt{nrow}}

Retorna o número de linhas:

\begin{lstlisting}
let n = count(df)           // total de observações
let n2 = nrow(df)           // alias de count(df)
let n3 = count(df, renda > 1000)   // observações onde renda > 1000
\end{lstlisting}

\section{Variáveis especiais}

Dentro de comandos que operam linha a linha (\lstinline{generate}, \lstinline{mutate},
\lstinline{filter}), duas variáveis implícitas estão disponíveis:

\begin{description}
  \item[\lstinline{_n}] Número da linha atual (começa em 1).
  \item[\lstinline{_N}] Total de linhas do DataFrame.
\end{description}

\begin{lstlisting}
generate df id = _n          // 1, 2, 3, ..., N
generate df peso = 1.0/_N    // peso uniforme
\end{lstlisting}

\section{Operadores de séries temporais}

Dentro de expressões de \lstinline{generate} (forma imperativa), operadores especiais
para séries temporais estão disponíveis:

\begin{center}
\begin{tabular}{cll}
\toprule
\textbf{Operador} & \textbf{Significado} & \textbf{Exemplo} \\
\midrule
\texttt{L.x}   & Lag (valor anterior)     & \texttt{L.pib}  → pib$_{t-1}$ \\
\texttt{F.x}   & Lead (valor seguinte)    & \texttt{F.taxa} → taxa$_{t+1}$ \\
\texttt{D.x}   & Diferença primeira       & \texttt{D.pib}  → pib$_t$ - pib$_{t-1}$ \\
\texttt{L2.x}  & Lag de ordem 2           & \texttt{L2.pib} → pib$_{t-2}$ \\
\bottomrule
\end{tabular}
\end{center}

\begin{lstlisting}[caption={Operadores de séries temporais}]
generate df dpib = D.pib      // variação do PIB
generate df pib_lag = L.pib   // PIB defasado um período
\end{lstlisting}
